\chapter{Appendix: Sums, Complex Numbers, and Calculus}

Following Appendix~A's foundation of sets and fields, this appendix gathers three
computational tools that appear throughout linear algebra: the \emph{summation sigma}
(for compact matrix and dot products), \emph{complex numbers} (essential for eigenvalues
and rotations), and a thin slice of \emph{calculus} (derivatives, gradients, and
integration for optimization and differential equations). None of it is difficult. But
if any of it feels shaky, the linear algebra built on top of it will feel shaky too, so we
make each piece concrete.

You do not need to read this appendix front to back. Treat it as a workshop to visit
whenever a main chapter writes something like $\sum_{k} a_{ik} b_{kj}$, or $e^{i\theta}$,
or $\nabla f$, and you want the exact definition of that symbol.

\begin{strangbox}
Strang's \emph{Introduction to Linear Algebra} assumes you already carry this toolbox: it writes
matrix products as sums over a shared index, jumps to complex eigenvalues for rotations without
pausing, and uses $\nabla$ and the matrix exponential when applications call for them. If those steps
ever seemed unexplained, read the matching section here \emph{before} the corresponding Strang
chapter, then return. \emph{After} Strang, this appendix is the place to re-anchor the three habits
that the whole subject rests on: a sum is a linear loop, complex multiplication is rotate-and-scale,
and setting a derivative to zero is how optimization becomes a linear system.
\end{strangbox}

\section{Sigma Notation: A Loop on Paper}
\label{ch39:sigma}

Suppose you want to add up the first five squares: $1 + 4 + 9 + 16 + 25$. Writing the whole sum is
fine for five terms. For five hundred it is hopeless, and for ``the first $n$'' it is impossible,
because you do not know $n$ yet. Mathematicians solved this with one symbol, the capital Greek
letter sigma, $\sum$, which means \emph{add these up}.

\begin{intuitionbox}
Read $\displaystyle\sum_{k=1}^{n} a_k$ out loud as a little program: ``start a running total at zero;
let $k$ count from $1$ up to $n$; each step, add $a_k$ to the total; when $k$ passes $n$, stop and
report the total.'' The $\sum$ is a \texttt{for} loop, and $a_k$ is the body of the loop. That is the
whole idea.
\end{intuitionbox}

The anatomy of the symbol:
\[
\sum_{\underbrace{k}_{\text{index}}=\underbrace{1}_{\text{start}}}^{\overbrace{n}^{\text{stop}}}
\underbrace{a_k}_{\text{term}}
\;=\; a_1 + a_2 + a_3 + \cdots + a_n.
\]
The letter $k$ is the \emph{index of summation}; it is a dummy variable, meaning it is a temporary
placeholder whose actual name has no meaning outside the loop. Writing $\sum_{k=1}^{n} a_k$ or
$\sum_{j=1}^{n} a_j$ or $\sum_{i=1}^{n} a_i$ all mean exactly the same number. The index exists only
inside the sum and never appears as a free variable in the final result.

\begin{example}[Unrolling a few sums]
\label{ch39:ex-unroll}
\[
\sum_{k=1}^{5} k^2 = 1^2 + 2^2 + 3^2 + 4^2 + 5^2 = 1+4+9+16+25 = 55.
\]
\[
\sum_{k=0}^{3} 2^k = 2^0 + 2^1 + 2^2 + 2^3 = 1 + 2 + 4 + 8 = 15.
\]
\[
\sum_{k=1}^{4} 3 = 3 + 3 + 3 + 3 = 12 \qquad(\text{the term does not have to mention } k).
\]
\end{example}

The last one is worth a pause: a sum of a constant $c$ taken $n$ times is just $nc$. In symbols
$\sum_{k=1}^{n} c = nc$. This is the first of several formulas we will collect.

\begin{recoverybox}
\textbf{Pause and think:} What is $\sum_{k=1}^5 10$? In your head: there is no $k$ in the summand, so you are simply adding five copies of $10$, giving $50$. What about $\sum_{k=1}^4 (k - (k-1))$? Before calculating anything, write the first two terms: $(1-0) + (2-1) + \cdots$. Every inner term cancels out, leaving only the very last upper value minus the starting lower value: $4 - 0 = 4$ (a telescoping sum).
\end{recoverybox}

\subsection{The three rules that make sigma algebra work}
\label{ch39:sigma-rules}

Because a sum is just repeated addition, it inherits the friendly properties of addition. Three of
them get used over and over, so we name them.

\begin{formalbox}
For any numbers $a_k, b_k$ and any constant $c$:
\begin{align}
\text{(Constant out)}\quad & \sum_{k=1}^{n} c\,a_k = c\sum_{k=1}^{n} a_k, \label{ch39:rule-const}\\[2pt]
\text{(Split a sum)}\quad & \sum_{k=1}^{n} (a_k + b_k) = \sum_{k=1}^{n} a_k + \sum_{k=1}^{n} b_k, \label{ch39:rule-split}\\[2pt]
\text{(Split the range)}\quad & \sum_{k=1}^{n} a_k = \sum_{k=1}^{m} a_k + \sum_{k=m+1}^{n} a_k \quad (1 \le m < n). \label{ch39:rule-range}
\end{align}
\end{formalbox}

Rule~\eqref{ch39:rule-const} is just the distributive law spread across many terms:
$ca_1 + ca_2 + \cdots = c(a_1 + a_2 + \cdots)$. Rule~\eqref{ch39:rule-split} is the commutativity of
addition: you may regroup $(a_1+b_1) + (a_2+b_2) + \cdots$ as $(a_1 + a_2 + \cdots) + (b_1 + b_2 +
\cdots)$. Rule~\eqref{ch39:rule-range} just says you may add the first chunk and the second chunk
separately. Together these three are \emph{linearity} of the sum, and linearity is the reason
$\sum$ plays so nicely with linear algebra.

\begin{trapbox}
A constant pulls out, but a \emph{product of two indexed things does not split}. In general
\[
\sum_k a_k b_k \;\ne\; \Bigl(\sum_k a_k\Bigr)\Bigl(\sum_k b_k\Bigr).
\]
Check it on $a = b = (1,1)$: the left side is $1\cdot1 + 1\cdot1 = 2$, the right side is
$(1+1)(1+1) = 4$. The left side $\sum_k a_k b_k$ is exactly the dot product $\vec{a}\cdot\vec{b}$,
and the whole point of the dot product is that it is \emph{not} the product of the sums. Keeping this
straight is what keeps matrix multiplication straight.
\end{trapbox}

\subsection{Four closed-form sums worth memorizing}
\label{ch39:closed-forms}

A few sums collapse to a clean formula. You will meet them when counting operations in elimination,
in variance formulas, and in the geometric series behind Markov chains and the Neumann series.

\begin{formalbox}
\begin{align}
\sum_{k=1}^{n} 1 &= n, \label{ch39:sum-ones}\\
\sum_{k=1}^{n} k &= \frac{n(n+1)}{2}, \label{ch39:sum-k}\\
\sum_{k=1}^{n} k^2 &= \frac{n(n+1)(2n+1)}{6}, \label{ch39:sum-ksq}\\
\sum_{k=0}^{n-1} r^k &= \frac{1 - r^n}{1 - r}\quad (r \ne 1), \qquad
\sum_{k=0}^{\infty} r^k = \frac{1}{1-r}\ \ (\abs{r}<1). \label{ch39:geometric}
\end{align}
\end{formalbox}

The triangular-number formula~\eqref{ch39:sum-k} has the famous one-line proof attributed to a young
Gauss: write the sum forwards and backwards and add columns,
\[
\begin{array}{rcccccc}
S &=& 1 &+& 2 &+& \cdots + n,\\
S &=& n &+& (n{-}1) &+& \cdots + 1,\\ \hline
2S &=& (n{+}1) &+& (n{+}1) &+& \cdots + (n{+}1),
\end{array}
\]
which is $n$ copies of $(n+1)$, so $2S = n(n+1)$ and $S = n(n+1)/2$. Figure~\ref{ch39:fig-gauss}
turns this into a picture: a staircase of dots, plus a second copy flipped on top, fills a clean
$n \times (n+1)$ rectangle.

\begin{figure}[H]
\centering
\begin{tikzpicture}[scale=0.42]
% staircase of blue dots: row i has i dots, i=1..5
\foreach \r in {1,...,5}{
  \foreach \c in {1,...,\r}{
    \fill[vecblue] (\c,\r) circle (0.28);
  }
}
% complementary staircase to fill the 5 x 6 rectangle, red
\foreach \r in {1,...,5}{
  \foreach \c in {1,...,6}{
    \pgfmathtruncatemacro{\lim}{\r}
    \ifnum\c>\lim
      \fill[vecred!75] (\c,\r) circle (0.28);
    \fi
  }
}
\draw[spacegray,thick] (0.5,0.5) rectangle (6.5,5.5);
\node[below] at (3.5,0.2) {\small width $n+1=6$};
\node[left] at (0.3,3) {\small height $n=5$};
\end{tikzpicture}
\caption{Gauss's sum as a picture, for $n=5$. The blue staircase is $1+2+3+4+5$. A red copy,
flipped, completes a $5\times 6$ rectangle of $30$ dots. So the blue half is $30/2 = 15 =
\tfrac{5\cdot 6}{2}$.}
\label{ch39:fig-gauss}
\end{figure}

Why does the sum of squares $\sum_{k=1}^n k^2 = \frac{n(n+1)(2n+1)}{6}$ scale as $n^3/3$? Think of
stacking discrete cubes to build a pyramid, or taking an integral: $\int_0^n x^2\,dx = n^3/3$.
Summing constants produces a linear result ($\sim n$), summing linear terms $k$ produces a
quadratic ($\sim n^2/2$), and summing squares must produce a cubic ($\sim n^3/3$). Expanding
$\frac{n(n+1)(2n+1)}{6} = \frac{n^3}{3} + \frac{n^2}{2} + \frac{n}{6}$ reveals that exact volume
scaling, with the smaller lower-degree terms accounting for the jagged edges of the discrete steps.

The geometric series~\eqref{ch39:geometric} is the most important of the four. Its derivation is a
telescoping trick. Let $S = 1 + r + r^2 + \cdots + r^{n-1}$. Multiply by $r$ and subtract:
\[
S - rS = (1 + r + \cdots + r^{n-1}) - (r + r^2 + \cdots + r^n) = 1 - r^n,
\]
because every middle term cancels. So $S(1-r) = 1 - r^n$, giving $S = (1-r^n)/(1-r)$. If
$\abs{r} < 1$ then $r^n \to 0$ as $n \to \infty$, and the infinite sum is $1/(1-r)$.

\begin{connectionbox}
The matrix version of $\dfrac{1}{1-r} = 1 + r + r^2 + \cdots$ is the \emph{Neumann series}: if a
square matrix $\mat{A}$ is small enough (its eigenvalues all have $\abs{\lambda} < 1$), then
\[
(\mat{I} - \mat{A})\inv = \mat{I} + \mat{A} + \mat{A}^2 + \mat{A}^3 + \cdots.
\]
The same telescoping proof works: $(\mat{I}-\mat{A})(\mat{I} + \mat{A} + \cdots + \mat{A}^{n-1}) =
\mat{I} - \mat{A}^n$, and $\mat{A}^n \to 0$. This is exactly the algebra behind ranking the web
(PageRank) and behind the steady state of a Markov chain. Scalar intuition transfers almost word for
word once you are willing to multiply matrices.
\end{connectionbox}

\subsection{Double sums and why matrix multiplication looks the way it does}
\label{ch39:double-sums}

When you nest one loop inside another you get a \emph{double sum}:
\[
\sum_{i=1}^{m}\sum_{j=1}^{n} a_{ij} \;=\; \sum_{i=1}^{m}\Bigl(\,\sum_{j=1}^{n} a_{ij}\Bigr).
\]
The inner sum (over $j$) runs to completion for each fixed value of the outer index $i$; then the
outer sum adds those partial totals. If you imagine the numbers $a_{ij}$ laid out in a grid with $i$
choosing the row and $j$ the column, a double sum just adds up every cell of the grid. And ``add up every cell'' does not care whether you sweep row-by-row or column-by-column, so you
may swap the order whenever the index ranges are independent:
\[
\sum_{i}\sum_{j} a_{ij} \;=\; \sum_{j}\sum_{i} a_{ij}.
\]
Figure~\ref{ch39:fig-doublesum} makes this physical: row totals added together or column totals added
together both reach the exact same grand sum.

\begin{figure}[H]
\centering
\begin{tikzpicture}[scale=0.95,>=Stealth]
  \begin{scope}[shift={(0,0)}]
    \node[anchor=south,font=\small\bfseries,text=spacegray!90!black] at (1.4,2.5) {(a) Sum rows first};
    \draw[fill=codebg,draw=spacegray!30,rounded corners=4pt] (-0.1,-0.3) rectangle (2.9,2.3);
    \foreach \r/\i in {1.8/1, 1.0/2, 0.2/3} {
      \foreach \c/\j in {0.4/1, 1.4/2, 2.4/3} {
        \node[font=\footnotesize] at (\c,\r) {$a_{\i\j}$};
      }
      \draw[->,vecblue,thick] (3.0,\r) -- (3.4,\r);
    }
    \node[right,vecblue,font=\scriptsize] at (3.4,1.8) {$r_1$};
    \node[right,vecblue,font=\scriptsize] at (3.4,1.0) {$r_2$};
    \node[right,vecblue,font=\scriptsize] at (3.4,0.2) {$r_3$};
    \draw[->,vecred,very thick] (4.0,1.5) -- (4.0,0.1) node[below,vecred,font=\footnotesize] {Total};
  \end{scope}

  \begin{scope}[shift={(6.0,0)}]
    \node[anchor=south,font=\small\bfseries,text=spacegray!90!black] at (1.4,2.5) {(b) Sum columns first};
    \draw[fill=codebg,draw=spacegray!30,rounded corners=4pt] (-0.1,-0.3) rectangle (2.9,2.3);
    \foreach \r/\i in {1.8/1, 1.0/2, 0.2/3} {
      \foreach \c/\j in {0.4/1, 1.4/2, 2.4/3} {
        \node[font=\footnotesize] at (\c,\r) {$a_{\i\j}$};
      }
    }
    \foreach \c/\j in {0.4/1, 1.4/2, 2.4/3} {
      \draw[->,targetgreen,thick] (\c,-0.4) -- (\c,-0.75);
    }
    \node[below,targetgreen,font=\scriptsize] at (0.4,-0.75) {$c_1$};
    \node[below,targetgreen,font=\scriptsize] at (1.4,-0.75) {$c_2$};
    \node[below,targetgreen,font=\scriptsize] at (2.4,-0.75) {$c_3$};
    \draw[->,vecred,very thick] (2.9,-1.05) -- (3.5,-1.05) node[right,vecred,font=\footnotesize] {Total};
  \end{scope}
\end{tikzpicture}
\caption{Swapping summation order in a double sum. In (a), the inner sum sweeps across each row to
produce row totals $r_i = \sum_j a_{ij}$, which the outer sum adds up. In (b), the inner sum sweeps
down each column to produce column totals $c_j = \sum_i a_{ij}$, which the outer sum combines.
Because every entry in the grid is counted exactly once with no omissions or duplicates, both paths
yield the identical grand total: $\sum_i \sum_j a_{ij} = \sum_j \sum_i a_{ij}$.}
\label{ch39:fig-doublesum}
\end{figure}

This single fact is the engine room of matrix algebra. The $(i,j)$ entry of a product $\mat{A}\mat{B}$
is a sum over a shared index:
\[
(\mat{A}\mat{B})_{ij} = \sum_{k=1}^{n} a_{ik} b_{kj}.
\]
Read it as the dot product of row $i$ of $\mat{A}$ with column $j$ of $\mat{B}$: $k$ marches along
the row and down the column simultaneously. The trace is a single sum down the diagonal,
$\tr(\mat{A}) = \sum_{i} a_{ii}$. For a product $\mat{A}\mat{B}$, the diagonal entries are
$(\mat{A}\mat{B})_{ii} = \sum_k a_{ik} b_{ki}$, so
\[
\tr(\mat{A}\mat{B}) = \sum_{i=1}^m \sum_{k=1}^n a_{ik} b_{ki}.
\]
Inside the double sum, $a_{ik}$ and $b_{ki}$ are ordinary scalar numbers, so their product
commutes: $a_{ik} b_{ki} = b_{ki} a_{ik}$. Swapping the order of summation now gives
\[
\tr(\mat{A}\mat{B}) = \sum_{i=1}^m \sum_{k=1}^n a_{ik} b_{ki}
= \sum_{k=1}^n \sum_{i=1}^m b_{ki} a_{ik}
= \sum_{k=1}^n (\mat{B}\mat{A})_{kk} = \tr(\mat{B}\mat{A}).
\]
The cyclic invariance of the trace is an immediate consequence of swapping two nested loops.

\begin{focusbox}
Hold onto three readings of $\sum$: (1) it is a \texttt{for} loop; (2) it is \emph{linear}, so
constants come out and sums split; (3) when two of them nest, you may swap the order. Almost every
``index gymnastics'' computation in this book is one of these three moves applied carefully.
\end{focusbox}

\begin{walkbox}
On your walk, write the $2\times 2$ product $\mat{A}\mat{B}$ entry by entry using only the formula
$(\mat{A}\mat{B})_{ij} = \sum_k a_{ik}b_{kj}$, with no matrix picture in front of you. Then check it
against the rows-times-columns picture you already know. Convince yourself the symbol and the picture
are the same object seen from two sides.
\end{walkbox}

\section{Complex Numbers: Rotation Made into Arithmetic}
\label{ch39:complex}

The eigenvalue chapter hits a wall the first time it meets a rotation: a $90^\circ$ turn in the plane
has \emph{no} real eigenvalue, because no real vector keeps its direction under a quarter turn. The
fix is to enlarge the number system. Complex numbers are what you get when you decide that the
equation $x^2 = -1$ should have a solution after all. What you gain is large: complex multiplication
\emph{is} rotation-and-scaling, which is precisely the geometry eigenvalues are trying to describe.

\subsection{The number $i$ and the complex plane}
\label{ch39:i}

Define a new number $i$ by the single rule
\[
i^2 = -1.
\]
A \emph{complex number} is then anything of the form $z = a + bi$ with $a, b$ real. We call $a =
\operatorname{Re}(z)$ the real part and $b = \operatorname{Im}(z)$ the imaginary part (note: $b$ is a
real number; ``imaginary'' is just a historical label). The set of all of them is $\C$.

The picture carries the idea. Plot $z = a + bi$ as the point $(a,b)$ in a plane, with the real
axis horizontal and the imaginary axis vertical. This is the \emph{complex plane}: geometrically it is
$\R^2$ equipped with complex multiplication.

\begin{historybox}
In 1545, the Italian physician and mathematician Gerolamo Cardano published \emph{Ars Magna} (The Great Art), presenting algebraic formulas for solving cubic ($x^3 + px = q$) and quartic polynomial equations. These formulas had been discovered through secret mathematical duels by Scipione del Ferro, Niccol\`o Tartaglia, and Cardano's student Lodovico Ferrari. But the formula ran headlong into an alarming mystery known as the \emph{casus irreducibilis} (the ``irreducible case''): whenever a cubic equation had three distinct real roots, the algebraic formula stubbornly demanded taking square roots of negative numbers. For instance, when applied to $x^3 = 15x + 4$, Cardano's formula produced
\[
  x = \sqrt[3]{2 + \sqrt{-121}} + \sqrt[3]{2 - \sqrt{-121}}.
\]
At the time, mathematicians possessed only geometric algebra inherited from Euclid and arithmetic on positive quantities; negative numbers were barely trusted, and the square root of a negative quantity was regarded as pure nonsense. Yet $x = 4$ was an obvious real solution ($4^3 = 64 = 15(4) + 4$)! Cardano was baffled, calling these expressions ``as subtle as they are useless'' and dismissing them as ``mental torture.''

In 1572, Rafael Bombelli of Bologna resolved the controversy in his treatise \emph{L'Algebra}. Bombelli had the insight to treat $\sqrt{-1}$ as an authentic quantity with consistent multiplication rules: $(+i)(+i) = -1$, $(-i)(-i) = -1$, and $(+i)(-i) = 1$. He discovered that $(2 + i)^3 = 2 + 11i = 2 + \sqrt{-121}$ and $(2 - i)^3 = 2 - 11i = 2 - \sqrt{-121}$. When added, the imaginary terms cancel out completely, leaving $(2+i) + (2-i) = 4$.

Suspicion lingered for two more centuries. Ren\'e Descartes coined the disparaging label ``imaginary'' in 1637 to imply these numbers were fictitious non-entities, and Gottfried Wilhelm Leibniz described $\sqrt{-1}$ as ``an amphibian between being and non-being.'' The fog lifted when Leonhard Euler introduced $e^{i\theta} = \cos\theta + i\sin\theta$ in 1748, and Caspar Wessel, Jean-Robert Argand, and Carl Friedrich Gauss gave complex numbers a physical home: points $(a,b)$ in a two-dimensional geometric plane where multiplication means rigid rotation and scaling.

This geometric view supported rigorous proofs of the Fundamental Theorem of Algebra (every $n$-th degree polynomial has exactly $n$ complex roots), electrical circuit analysis, and quantum mechanics, where physical states live in complex vector spaces. Yet deep mysteries remain open: Bernhard Riemann's 1859 hypothesis asserts that all non-trivial zeros of the Riemann zeta function $\zeta(s) = \sum_{n=1}^\infty n^{-s}$ (a function built by summing reciprocal powers across the complex plane) lie precisely on the ``critical line'' $\operatorname{Re}(s) = \tfrac12$. Intimately connected to the distribution of prime numbers and the eigenvalues of chaotic quantum systems, proving or disproving the Riemann Hypothesis remains the most celebrated open problem in all of mathematics.
\end{historybox}

\begin{figure}[H]
\centering
\begin{tikzpicture}[scale=1.05,>=Stealth]
\draw[->,spacegray] (-0.6,0) -- (3.6,0) node[right] {\small real axis};
\draw[->,spacegray] (0,-2.5) -- (0,2.6) node[above] {\small imaginary axis};
\foreach \x in {1,2}\draw[spacegray] (\x,0.06)--(\x,-0.06) node[below,font=\tiny]{\x};
\draw[spacegray] (3,0.06)--(3,-0.06) node[below left,font=\tiny,inner sep=0.5pt]{3};
\draw[spacegray] (0.06,1)--(-0.06,1) node[left,font=\tiny]{$i$};
\draw[spacegray] (0.06,2)--(-0.06,2) node[left,font=\tiny]{$2i$};
\draw[spacegray] (0.06,-1)--(-0.06,-1) node[left,font=\tiny]{$-i$};
\draw[spacegray] (0.06,-2)--(-0.06,-2) node[left,font=\tiny]{$-2i$};
% z = 3 + 2i
\draw[->,vecblue,very thick] (0,0) -- (3,2);
\fill[vecblue] (3,2) circle (1.6pt);
\node[vecblue,above right] at (3,2) {\small $z=3+2i$};
\node[vecblue,above left] at (1.5,1.0) {\small $r=\abs{z}$};
% angle theta
\draw[targetgreen] (0.7,0) arc (0:33.69:0.7);
\node[targetgreen,right] at (0.72,0.22) {\small $\theta$};
% conjugate
\draw[->,structurepurple,thick] (0,0) -- (3,-2);
\fill[structurepurple] (3,-2) circle (1.6pt);
\node[structurepurple,below right] at (3,-2) {\small $\bar z = 3-2i$};
% vertical drop connecting conjugates
\draw[spacegray,dashed] (3,2) -- (3,-2);
\path (0,-2.6) (3.6,2.7);
\end{tikzpicture}
\caption{A complex number $z = a+bi$ lives at the point $(a,b)$ in the complex plane. Its length
$r = \abs{z}$ and angle $\theta$ from the positive real axis are its polar coordinates. The
conjugate $\bar z = a - bi$ is its mirror image across the real axis; the vertical dashed line
emphasizes that both share the same real component $a=3$.}
\label{ch39:fig-plane}
\end{figure}

\subsection{Arithmetic: add like vectors, multiply by rotating}
\label{ch39:complex-arith}

\textbf{Addition} is componentwise, exactly like adding vectors in $\R^2$:
\[
(a+bi) + (c+di) = (a+c) + (b+d)i.
\]
\textbf{Multiplication} uses ordinary distribution, then collapses $i^2$ to $-1$:
\[
(a+bi)(c+di) = ac + adi + bci + bd\,i^2 = (ac - bd) + (ad + bc)i.
\]

\begin{example}[A product worked term by term]
\label{ch39:ex-mult}
\[
(2 + 3i)(1 - 4i) = 2 - 8i + 3i - 12 i^2 = 2 - 8i + 3i + 12 = 14 - 5i.
\]
The $-12 i^2$ became $+12$ because $i^2 = -1$. This sign change is the defining feature
of complex multiplication.
\end{example}

Two more pieces of vocabulary. The \emph{conjugate} of $z = a + bi$ is $\bar z = a - bi$, the
reflection of $z$ across the real axis. The \emph{modulus} (or absolute value, or length) is
\[
\abs{z} = \sqrt{a^2 + b^2},
\]
the ordinary distance from the origin. Conjugate and modulus are tied together by a clean identity:
\[
z\bar z = (a+bi)(a-bi) = a^2 - (bi)^2 = a^2 + b^2 = \abs{z}^2.
\]
So $z\bar z$ is always a nonnegative \emph{real} number, namely $\abs{z}^2$. This is exactly how you
\emph{divide}: to compute $1/z$, multiply top and bottom by $\bar z$ to make the denominator real.

\begin{example}[Division by realizing the denominator]
\label{ch39:ex-div}
\[
\frac{1}{2+3i} = \frac{1}{2+3i}\cdot\frac{2-3i}{2-3i}
= \frac{2-3i}{2^2+3^2} = \frac{2-3i}{13} = \frac{2}{13} - \frac{3}{13}i.
\]
\end{example}

\begin{recoverybox}
\textbf{Pause and think:} What is $i^4$? In your head: $i^2 = -1$, so $i^4 = (i^2)^2 = (-1)^2 = 1$.
Geometrically, multiplying by $i$ is a $90^\circ$ quarter-turn; four quarter-turns make a complete
$360^\circ$ circle, returning right back to $+1$. What is the modulus of $3 + 4i$? By Pythagoras,
$\abs{3+4i} = \sqrt{3^2 + 4^2} = \sqrt{25} = 5$.
\end{recoverybox}

\begin{trapbox}
The modulus is \emph{not} additive: $\abs{z+w} \ne \abs{z} + \abs{w}$ in general (the triangle
inequality only gives $\le$). But the modulus \emph{is} multiplicative: $\abs{zw} = \abs{z}\,\abs{w}$
always. Multiplication respects lengths; addition does not. The next subsection explains why, in one
picture.
\end{trapbox}

\subsection{Polar form and Euler's formula}
\label{ch39:euler}

Switch a complex number to polar coordinates. If $z$ has length $r = \abs{z}$ and makes angle
$\theta$ with the positive real axis, then $a = r\cos\theta$ and $b = r\sin\theta$, so
\[
z = r(\cos\theta + i\sin\theta).
\]
Now multiply two numbers written this way, $z_1 = r_1(\cos\alpha + i\sin\alpha)$ and $z_2 =
r_2(\cos\beta + i\sin\beta)$. Expand and use the angle-addition formulas for sine and cosine:
\[
z_1 z_2 = r_1 r_2\bigl[(\cos\alpha\cos\beta - \sin\alpha\sin\beta) + i(\sin\alpha\cos\beta +
\cos\alpha\sin\beta)\bigr] = r_1 r_2\bigl(\cos(\alpha+\beta) + i\sin(\alpha+\beta)\bigr).
\]

\begin{intuitionbox}
\textbf{To multiply complex numbers, multiply the lengths and add the angles.} That is the entire
geometric content. Multiplication by a complex number is a \emph{rotate-and-scale} of the plane.
Multiplying by $i$ (length $1$, angle $90^\circ$) is a pure quarter-turn; check it: $i\cdot 1 = i$
(the point $1$ swings up to $i$), and $i\cdot i = -1$ (another quarter-turn lands on $-1$). The
mysterious rule $i^2 = -1$ is just ``two quarter-turns make a half-turn.''
\end{intuitionbox}

Figure~\ref{ch39:fig-mult} shows the rule in one picture: stack the two angles tip to tail and
multiply the two lengths.

\begin{figure}[H]
\centering
\begin{tikzpicture}[scale=1.25,>=Stealth]
\draw[->,spacegray] (-0.5,0) -- (2.7,0) node[right] {\scriptsize Re};
\draw[->,spacegray] (0,-0.5) -- (0,2.3) node[above] {\scriptsize Im};
% z1: r1=1.2, alpha=25 deg
\def\rone{1.2}\def\al{25}
% z2: r2=1.4, beta=40 deg
\def\rtwo{1.4}\def\be{40}
% product: r=1.68, angle=65
\pgfmathsetmacro{\rp}{\rone*\rtwo}
\pgfmathsetmacro{\ab}{\al+\be}
\draw[->,vecblue,very thick] (0,0) -- (\al:\rone);
\node[vecblue,below right,font=\tiny] at (\al:\rone) {$z_1=r_1e^{i\alpha}$};
\draw[->,targetgreen,very thick] (0,0) -- (\be:\rtwo);
\node[targetgreen,right,font=\tiny] at (\be:\rtwo) {$z_2=r_2e^{i\beta}$};
\draw[->,vecred,very thick] (0,0) -- (\ab:\rp);
\node[vecred,above right,font=\tiny] at (\ab:\rp) {$z_1z_2=r_1r_2\,e^{i(\alpha+\beta)}$};
% angle arcs
\draw[vecblue] (0.45,0) arc (0:\al:0.45);
\node[vecblue,font=\tiny] at (12.5:0.58) {$\alpha$};
\draw[targetgreen,dashed] (\al:0.75) arc (\al:\ab:0.75);
\node[targetgreen,font=\tiny,fill=white,inner sep=1pt] at (55:0.85) {$\beta$};
\end{tikzpicture}
\caption{Multiplication adds angles and multiplies lengths. The product $z_1 z_2$ (red) sits at angle
$\alpha+\beta$ with length $r_1 r_2$. Here the green dashed arc shows the angle $\beta$ stacked
directly onto $\alpha$, and the red arrow is longer than either factor because $r_1 r_2 > r_1, r_2$
when both lengths exceed $1$.}
\label{ch39:fig-mult}
\end{figure}

The ``add the angles'' rule is the exact rule for exponents, $e^\alpha e^\beta = e^{\alpha+\beta}$,
and Euler made that connection precise.

\begin{formalbox}[Euler's formula]
For any real $\theta$,
\[
e^{i\theta} = \cos\theta + i\sin\theta.
\]
Consequently every complex number has the compact polar form $z = r\,e^{i\theta}$, with $r =
\abs{z}$ and $\theta$ its angle. In this notation, multiplication is just
$ (r_1 e^{i\alpha})(r_2 e^{i\beta}) = r_1 r_2\, e^{i(\alpha+\beta)}$, and the special value at
$\theta = \pi$ is the celebrated identity $e^{i\pi} + 1 = 0$.
\end{formalbox}

Why is Euler's formula true? The quickest rigorous argument comes from power series. The exponential,
sine, and cosine each have a Taylor series (we meet Taylor series in Section~\ref{ch39:calculus}):
\[
e^x = \sum_{k=0}^{\infty}\frac{x^k}{k!},\qquad
\cos\theta = \sum_{k=0}^{\infty}\frac{(-1)^k\theta^{2k}}{(2k)!},\qquad
\sin\theta = \sum_{k=0}^{\infty}\frac{(-1)^k\theta^{2k+1}}{(2k+1)!}.
\]
Substitute $x = i\theta$ into the exponential series and use the cycle $i^0=1,\ i^1=i,\ i^2=-1,\
i^3=-i,\ i^4 = 1,\dots$:
\[
e^{i\theta} = 1 + i\theta + \frac{(i\theta)^2}{2!} + \frac{(i\theta)^3}{3!} + \cdots
= \underbrace{\Bigl(1 - \tfrac{\theta^2}{2!} + \tfrac{\theta^4}{4!} - \cdots\Bigr)}_{\cos\theta}
+ i\underbrace{\Bigl(\theta - \tfrac{\theta^3}{3!} + \tfrac{\theta^5}{5!} - \cdots\Bigr)}_{\sin\theta}.
\]
The real terms reassemble the cosine series and the imaginary terms reassemble the sine series. Done.

\begin{figure}[H]
\centering
\begin{tikzpicture}[scale=1.5,>=Stealth]
\draw[->,spacegray] (-1.35,0) -- (1.4,0) node[right] {\scriptsize Re};
\draw[->,spacegray] (0,-1.35) -- (0,1.4) node[above] {\scriptsize Im};
\draw[spacegray,thin] (0,0) circle (1);
\node[spacegray,above left,font=\tiny] at (-0.72,0.72) {unit circle};
% angle theta = 50 deg point
\def\ang{50}
\draw[->,vecblue,very thick] (0,0) -- (\ang:1);
\fill[vecblue] (\ang:1) circle (0.8pt);
\node[vecblue,above right] at (\ang:1) {\scriptsize $e^{i\theta}=\cos\theta+i\sin\theta$};
% horizontal shadow (cos) along the real axis
\draw[targetgreen,very thick] (0,0) -- ({cos(\ang)},0);
\node[targetgreen,below,font=\tiny] at ({0.5*cos(\ang)},-0.04) {$\cos\theta$};
% vertical shadow (sin), dropped from the point to the real axis
\draw[vecred,very thick] ({cos(\ang)},0) -- ({cos(\ang)},{sin(\ang)});
\node[vecred,right,font=\tiny] at ({cos(\ang)},{0.5*sin(\ang)}) {$\sin\theta$};
\draw[orange] (0.32,0) arc (0:\ang:0.32);
\node[orange,font=\tiny] at (0.5,0.18) {$\theta$};
\end{tikzpicture}
\caption{Euler's formula geometrically: $e^{i\theta}$ rides around the unit circle. Its horizontal
shadow is $\cos\theta$, its vertical shadow is $\sin\theta$. Increasing $\theta$ just walks the point
counterclockwise; multiplying by $e^{i\theta}$ rotates the whole plane by $\theta$.}
\label{ch39:fig-euler}
\end{figure}

\begin{example}[De Moivre powers]
\label{ch39:ex-demoivre}
Polar form makes powers trivial. Since $(re^{i\theta})^n = r^n e^{in\theta}$, raising to a power
scales the length and multiplies the angle. Take $z = 1 + i$, which has $r = \sqrt2$ and $\theta =
45^\circ = \pi/4$. Then
\[
(1+i)^8 = (\sqrt2)^8\, e^{i\cdot 8\pi/4} = 16\, e^{i 2\pi} = 16.
\]
Try multiplying $(1+i)$ by itself eight times by hand to appreciate how much labor polar form just
saved.
\end{example}

\begin{applicationbox}
The discrete Fourier transform, which underlies JPEG, MP3, and every digital signal pipeline, is
built entirely from the numbers $\omega^k = e^{-2\pi i k/n}$ --- equally spaced points around the unit
circle, called \emph{roots of unity} because their $n$-th power equals $1$. The transform writes a
signal as a combination of these pure rotations at different speeds. Because each $\omega^k$ has
modulus $1$, the transform never inflates energy; that length-preservation is exactly the statement
that the Fourier matrix is \emph{unitary}, the complex cousin of an orthogonal matrix (its columns
are perpendicular unit vectors in $\C^n$, preserving lengths and inner products during transformation).
The Fourier chapter of this book takes $e^{i\theta}$ as its alphabet, so being fluent here pays off
directly there.
\end{applicationbox}

\begin{connectionbox}
Complex numbers are why the eigenvalues of a real matrix can come in conjugate pairs $\lambda =
a \pm bi$. Geometrically such a pair means the matrix is doing a rotation-and-scale inside some
two-dimensional invariant plane: the modulus $\abs{\lambda} = \sqrt{a^2+b^2}$ is the per-step
\emph{stretch}, and the argument $\theta = \arg\lambda$ (the angle the number makes with the
positive real axis) is the per-step \emph{turn}. A discrete dynamical system $\vec{x}_{k+1} =
\mat{A}\vec{x}_k$ spirals outward when $\abs{\lambda} > 1$, spirals inward when $\abs{\lambda} < 1$,
and circles forever when $\abs{\lambda} = 1$. The whole qualitative story lives in one complex number.
\end{connectionbox}

\begin{labbox}
Verify the rotation claim with a $2\times 2$ matrix. The rotation by $90^\circ$ is
$\mat{R} = \begin{bmatrix} 0 & -1\\ 1 & 0\end{bmatrix}$. Compute its eigenvalues by solving
$\det(\mat{R} - \lambda \mat{I}) = \lambda^2 + 1 = 0$, which gives $\lambda = \pm i$. There are no
real eigenvectors (no real arrow survives a quarter-turn pointing the same way), and the eigenvalues
are pure imaginary with $\abs{\lambda} = 1$ --- a turn with no stretch, exactly as the picture
demands. Now do the same for $\begin{bmatrix} 0 & -2\\ 2 & 0\end{bmatrix}$ and watch $\abs{\lambda}$
become $2$.
\end{labbox}

\section{Just Enough Calculus}
\label{ch39:calculus}

Two chapters of this book reach outside linear algebra into calculus. The differential-equations
chapter needs \emph{derivatives} (to say what ``rate of change'' means) and a little
\emph{integration} (to define the matrix exponential and to take inner products of functions). The
optimization chapter needs the \emph{gradient} (the multivariable derivative) to find where a
function is flat, which is where minima live. This section builds exactly those tools and stops.

\subsection{The derivative as a slope, and as a best linear fit}
\label{ch39:derivative}

Start with one variable. The derivative of $f$ at a point $x$ measures how fast $f$ is changing
there. Geometrically it is the slope of the tangent line: zoom in on the graph near $x$ until the
curve looks straight, and the derivative is that straight line's slope.

\begin{formalbox}
The \emph{derivative} of $f$ at $x$ is the limit of slopes of secant lines (straight lines connecting
two nearby points on the curve),
\[
f'(x) = \frac{df}{dx} = \lim_{h\to 0}\frac{f(x+h) - f(x)}{h},
\]
when the limit exists. Equivalently, near $x$ the function is well approximated by a line:
$f(x+h) \approx f(x) + f'(x)\,h$ for small $h$.
\end{formalbox}

\begin{figure}[H]
\centering
\begin{tikzpicture}[scale=1.0,>=Stealth]
\begin{axis}[
  width=8.8cm, height=5.4cm,
  axis lines=middle, xlabel={\small $x$}, ylabel={\small $f(x)$},
  xtick=\empty, ytick=\empty, clip=false,
  xmin=-0.3,xmax=3.2, ymin=-0.3, ymax=4.6,
  domain=0:3, samples=80]
\addplot[vecblue,very thick] {0.5*x^2};
% tangent at x=1.6: f=0.5*2.56=1.28, f'=x=1.6
\addplot[vecred,thick,domain=0.4:2.8] {1.28 + 1.6*(x-1.6)};
% secant from x=1.6 to x=2.6
\addplot[targetgreen,dashed,thick,domain=0.9:2.9] {1.28 + ((0.5*2.6^2-1.28)/1.0)*(x-1.6)};
\addplot[only marks,mark=*,vecblue,mark size=1.4pt] coordinates {(1.6,1.28)};
\addplot[only marks,mark=*,targetgreen,mark size=1.4pt] coordinates {(2.6,3.38)};
\node[vecred,anchor=west,font=\small] at (axis cs:2.0,1.05) {tangent};
\node[targetgreen,anchor=south east,font=\small] at (axis cs:2.45,3.5) {secant};
\end{axis}
\end{tikzpicture}
\caption{The secant slope $\frac{f(x+h)-f(x)}{h}$ (green) is the average rate of change over a step
$h$. As the step shrinks, the secant pivots toward the tangent (red), whose slope is the derivative
$f'(x)$.}
\label{ch39:fig-derivative}
\end{figure}

\begin{example}[A derivative from the definition]
\label{ch39:ex-deriv}
Let $f(x) = x^2$. Then
\[
\frac{f(x+h)-f(x)}{h} = \frac{(x+h)^2 - x^2}{h} = \frac{2xh + h^2}{h} = 2x + h \xrightarrow[h\to 0]{} 2x.
\]
So $f'(x) = 2x$. The $h$ in the numerator is what lets the fraction survive division by $h$; sending
$h\to 0$ at the end discards the leftover $h$ and leaves the clean slope $2x$.
\end{example}

You will not re-derive every derivative from the limit. A handful of rules and a short table cover
everything the main chapters use.

\begin{formalbox}[The rules you actually need]
\setlength{\arraycolsep}{3pt}
\begin{align*}
&\text{Power:} && (x^n)' = n x^{n-1}, & &\text{Exponential:} && (e^{ax})' = a\,e^{ax},\\
&\text{Sine:} && (\sin x)' = \cos x, & &\text{Cosine:} && (\cos x)' = -\sin x,\\
&\text{Constant/linear:} && (c)' = 0,\ (x)' = 1, & &\text{Sum:} && (f+g)' = f' + g',\\
&\text{Constant multiple:} && (cf)' = c f', & &\text{Product:} && (fg)' = f'g + fg',\\
&\text{Chain:} && \bigl(f(g(x))\bigr)' = f'(g(x))\,g'(x). & & &&
\end{align*}
\end{formalbox}

The sum and constant-multiple rules say the derivative is a \emph{linear} operation: $(af + bg)' = af'
+ bg'$. That linearity is why differentiation itself can be treated as a matrix in the
change-of-basis and functional-analysis chapters --- a linear map on a space of functions. The chain
rule is the one to be most careful with; it forms the computational engine behind backpropagation.

\begin{example}[The chain rule in action]
\label{ch39:ex-chain}
Differentiate $f(x) = e^{-x^2}$ (the bell-curve shape). Outer function $e^{u}$ with $u = -x^2$:
\[
f'(x) = e^{-x^2}\cdot\frac{d}{dx}(-x^2) = e^{-x^2}\cdot(-2x) = -2x\,e^{-x^2}.
\]
This is zero only at $x = 0$, the peak of the bell, which matches the picture: the curve is flat
exactly at its top.
\end{example}

\begin{recoverybox}
\textbf{Pause and think:} For $f(x) = x^3$, what is the slope of the tangent line at $x=0$? By the
power rule, $f'(x) = 3x^2$, so $f'(0) = 0$: the cubic curve momentarily flattens out horizontal at
the origin before continuing its climb. What about $f(x) = 5x + 7$? Its derivative is $5$ everywhere,
confirming that the tangent slope of a straight line is always its constant rate of change.
\end{recoverybox}

\subsection{From slope to optimization: where the derivative is zero}
\label{ch39:optimization-1d}

The reason calculus shows up in the optimization chapter is one sentence: \emph{at a smooth minimum
or maximum, the slope is zero}. If $f'(x_\star) \ne 0$ then the tangent line tilts, and stepping
slightly downhill makes $f$ smaller, so $x_\star$ cannot be a minimum. Setting $f'(x) = 0$ and
solving locates the candidates.

\begin{example}[Least squares in one variable]
\label{ch39:ex-leastsq}
Suppose data points have heights $y_1, \dots, y_n$ and you want the single constant $c$ that best
fits them in the least-squares sense, i.e. the $c$ minimizing
\[
E(c) = \sum_{k=1}^{n} (y_k - c)^2.
\]
Differentiate with respect to $c$, using the chain rule on each term and the linearity of $\sum$:
\[
E'(c) = \sum_{k=1}^{n} 2(y_k - c)(-1) = -2\sum_{k=1}^{n}(y_k - c).
\]
Set $E'(c) = 0$: then $\sum_k (y_k - c) = 0$, i.e. $\sum_k y_k = nc$, so $c = \frac1n\sum_k y_k$.
The best constant fit is the \emph{average}. This little calculation is the seed of the whole
least-squares and regression story, where ``set the derivative to zero'' becomes the normal
equations $\mat{A}\T\mat{A}\hat{\vec{x}} = \mat{A}\T\vec{b}$.
\end{example}

\subsection{Gradients: the derivative for many variables}
\label{ch39:gradient}

When $f$ depends on several variables, $f(x_1, \dots, x_n)$, there is a slope in each coordinate
direction. The \emph{partial derivative} $\partial f/\partial x_j$ is just the ordinary derivative in
the $x_j$ direction, holding the other variables fixed. Stack all of them into a vector and you get
the \emph{gradient}.

\begin{formalbox}
The \emph{gradient} of $f\colon \R^n \to \R$ is the vector of its partial derivatives,
\[
\nabla f = \colvec{\partial f/\partial x_1 \\ \partial f/\partial x_2 \\ \vdots \\ \partial f/\partial x_n}.
\]
Near a point $\vec{x}$, the function is approximated by the linear formula $f(\vec{x} + \vec{h})
\approx f(\vec{x}) + \nabla f(\vec{x})\cdot \vec{h}$, the multivariable analog of $f(x+h)\approx f(x) +
f'(x)h$.
\end{formalbox}

The geometry of the gradient is worth holding onto.

\begin{intuitionbox}
At any point, $\nabla f$ points in the direction of \emph{steepest increase} of $f$, and its length
is how steep that climb is. Walk along $\nabla f$ and $f$ rises as fast as possible; walk along
$-\nabla f$ and $f$ falls as fast as possible (this is \emph{gradient descent}, an iterative
algorithm that repeatedly steps downhill to find a minimum); walk \emph{perpendicular} to $\nabla f$
and $f$ does not change at all (you are following a \emph{level curve}, a contour of constant height,
like an elevation line on a topographic map). A contour map of a mountain makes this visible: the
gradient is the arrow pointing straight uphill, always strictly at right angles to the contour lines.
\end{intuitionbox}

\begin{figure}[H]
\centering
\begin{tikzpicture}[scale=1.2,>=Stealth]
  \node[anchor=south,font=\footnotesize\scshape,text=spacegray!90!black] at (0,2.0) {Contours and Gradient Direction};
  \draw[fill=codebg,draw=spacegray!30,rounded corners=4pt] (-2.7,-1.8) rectangle (2.8,1.9);

  \draw[->,gray!50] (-2.3,0) -- (2.3,0) node[right,font=\scriptsize] {$x_1$};
  \draw[->,gray!50] (0,-1.5) -- (0,1.6) node[above,font=\scriptsize] {$x_2$};

  % Shaded level curves x1^2 + 2 x2^2 = c
  \draw[fill=sourceblue!6,draw=sourceblue!30] (0,0) ellipse (1.732 and 1.225);
  \draw[fill=sourceblue!14,draw=sourceblue!60,line width=1.1pt] (0,0) ellipse (1.311 and 0.927);
  \draw[fill=sourceblue!22,draw=sourceblue!90] (0,0) ellipse (0.707 and 0.5);

  % Evaluation point on middle contour: (1.0, 0.6)
  \coordinate (pt) at (1.0,0.6);
  \fill[black] (pt) circle (2pt) node[below left=1pt,font=\scriptsize]{$\vec{x}$};

  % Tangent line segment to the contour at (1.0, 0.6)
  \draw[structurepurple,thick,dashed] (1.58,0.12) node[above right=1pt,font=\tiny,structurepurple] {tangent}
    -- (0.42,1.08);

  % Gradient arrow (pointing uphill orthogonal to contour)
  \draw[->,vecred,line width=1.6pt] (pt) -- (1.65,1.38)
    node[above right=1pt,font=\small]{$\nabla f(\vec{x})$};

  % Right angle indicator between gradient and tangent
  \draw[spacegray!80,thick] (1.090,0.708) -- (0.982,0.798) -- (0.892,0.690);

  \node[draw=sourceblue!70,fill=white,rounded corners=2pt,font=\tiny] at (1.3,-1.42) {Level curve: $f(\vec{x})=c$};
  \node[draw=vecred!70,fill=white,rounded corners=2pt,font=\tiny] at (-1.3,-1.42) {$\nabla f \perp$ tangent line};
\end{tikzpicture}
\caption{Level curves of the quadratic bowl $f(x_1,x_2) = x_1^2 + 2x_2^2$. The gradient vector
$\nabla f = (2x_1, 4x_2)$ points in the direction of steepest uphill ascent and is strictly
perpendicular to the level contour through $\vec{x}$, as shown by the right-angle indicator against
the purple tangent line. At the minimum $(0,0)$, the gradient vanishes ($\nabla f = \vec{0}$).}
\label{ch39:fig-gradient}
\end{figure}

\begin{example}[Computing a gradient]
\label{ch39:ex-grad}
Let $f(x_1, x_2) = x_1^2 + 2x_2^2$. Holding $x_2$ fixed, $\partial f/\partial x_1 = 2x_1$; holding
$x_1$ fixed, $\partial f/\partial x_2 = 4x_2$. So
\[
\nabla f = \colvec{2x_1 \\ 4x_2}.
\]
The only point where $\nabla f = \vec{0}$ is the origin, and the bowl shape tells us it is a minimum.
\end{example}

\begin{example}[The gradient of a quadratic form]
\label{ch39:ex-quadform}
The optimization chapter lives on quadratic forms (functions where every term is degree two or one),
$f(\vec{x}) = \tfrac12 \vec{x}\T\mat{A}\vec{x} - \vec{b}\T\vec{x}$ with $\mat{A}$ symmetric. Its
gradient is the tidy
\[
\nabla f(\vec{x}) = \mat{A}\vec{x} - \vec{b},
\]
so setting $\nabla f = \vec{0}$ recovers the linear system $\mat{A}\vec{x} = \vec{b}$. When $\mat{A}$
is positive definite (all eigenvalues $\lambda_i > 0$), the surface curves upward in every direction
like a bowl, guaranteeing that this critical point is a unique global minimum. Let us verify the
$2\times2$ case with $\mat{A} = \begin{bmatrix}2 & 0\\ 0 & 4\end{bmatrix}$ and $\vec{b} = \vec{0}$:
then $\tfrac12\vec{x}\T\mat{A}\vec{x} = x_1^2 + 2x_2^2$, whose gradient we just found to be
$(2x_1, 4x_2) = \mat{A}\vec{x}$. The abstract formula and the hand computation agree.
\end{example}

\begin{formalbox}
\textbf{Matrix calculus derivatives, step by step.} Where does $\nabla_{\vec{x}}(\vec{x}\T\mat{A}\vec{x}) = 2\mat{A}\vec{x}$ come from? Expand the quadratic form as an index sum:
\[
  f(\vec{x}) = \vec{x}\T\mat{A}\vec{x} = \sum_{i=1}^n \sum_{j=1}^n A_{ij} x_i x_j.
\]
To find component $k$ of the gradient, differentiate with respect to $x_k$. The variable $x_k$ appears whenever $i=k$ or $j=k$:
\begin{itemize}
  \item When $i=k$ and $j \ne k$: the terms are $\sum_{j \ne k} A_{kj} x_k x_j$, with derivative $\sum_{j \ne k} A_{kj} x_j$.
  \item When $j=k$ and $i \ne k$: the terms are $\sum_{i \ne k} A_{ik} x_i x_k$, with derivative $\sum_{i \ne k} A_{ik} x_i$.
  \item When $i=j=k$: the single diagonal term is $A_{kk} x_k^2$, with derivative $2 A_{kk} x_k$.
\end{itemize}
Combining all three pieces:
\[
  \frac{\partial}{\partial x_k}(\vec{x}\T\mat{A}\vec{x})
  = \sum_{j=1}^n A_{kj} x_j + \sum_{i=1}^n A_{ik} x_i
  = (\mat{A}\vec{x})_k + (\mat{A}\T\vec{x})_k.
\]
Stacking all $k$ rows yields the general matrix calculus identity:
\[
  \boxed{\;\nabla_{\vec{x}}(\vec{x}\T\mat{A}\vec{x}) = (\mat{A} + \mat{A}\T)\vec{x}.\;}
\]
When $\mat{A}$ is symmetric ($\mat{A} = \mat{A}\T$), this becomes $2\mat{A}\vec{x}$. Dividing by $2$ gives $\nabla(\tfrac12\vec{x}\T\mat{A}\vec{x}) = \mat{A}\vec{x}$. Similarly, for a linear form $\vec{b}\T\vec{x} = \sum b_i x_i$, $\frac{\partial}{\partial x_k}(\vec{b}\T\vec{x}) = b_k$, so $\nabla_{\vec{x}}(\vec{b}\T\vec{x}) = \vec{b}$.
\end{formalbox}

\begin{applicationbox}
Training a neural network is gradient descent on a function of millions of variables. The loss $L$ is
a scalar measuring how wrong the network is; its gradient $\nabla L$ with respect to all the weights
is computed by \emph{backpropagation}, which is the chain rule from
Section~\ref{ch39:derivative} applied layer by layer. Each training step nudges the weights along
$-\nabla L$. Every matrix multiply inside a network is the
linear algebra; every gradient step is this calculus.
\end{applicationbox}

\subsection{A little integration}
\label{ch39:integration}

The differential-equations chapter needs integration in two small ways, so we collect just those.

First, the integral as \emph{accumulated area / undoing a derivative}. The definite integral
$\int_a^b f(x)\,dx$ is the signed area under the graph of $f$ between $a$ and $b$. The Fundamental
Theorem of Calculus (the foundational principle establishing that differentiation and integration are
inverse operations) ties it to derivatives: if $F' = f$ then
\[
\int_a^b f(x)\,dx = F(b) - F(a).
\]
Integration is the reverse of differentiation, just as $\sum$ is a discrete cousin of $\int$ (both
accumulate). For instance, since $(e^{ax})' = a e^{ax}$, we get $\int e^{ax}\,dx = \frac1a e^{ax} +
C$. That single antiderivative is what lets the matrix exponential $e^{\mat{A}t} = \sum_{k=0}^\infty
\frac{(\mat{A}t)^k}{k!}$ solve $\frac{d}{dt}\vec{x} = \mat{A}\vec{x}$: differentiate the series term by
term and the same series reappears, multiplied by $\mat{A}$.

\begin{example}[Differentiating the matrix-exponential series]
\label{ch39:ex-matexp}
Treat the series $e^{\mat{A}t} = \mat{I} + \mat{A}t + \frac{\mat{A}^2 t^2}{2!} + \frac{\mat{A}^3
t^3}{3!} + \cdots$ and differentiate term by term in $t$ using the power rule:
\[
\frac{d}{dt}e^{\mat{A}t} = \mat{A} + \mat{A}^2 t + \frac{\mat{A}^3 t^2}{2!} + \cdots
= \mat{A}\Bigl(\mat{I} + \mat{A}t + \frac{\mat{A}^2 t^2}{2!} + \cdots\Bigr) = \mat{A}\,e^{\mat{A}t}.
\]
So $\vec{x}(t) = e^{\mat{A}t}\vec{x}(0)$ satisfies $\vec{x}\,'(t) = \mat{A}\vec{x}(t)$. The scalar
fact $(e^{at})' = a e^{at}$ became its matrix twin with no new ideas, only the power rule and the
linearity of differentiation.
\end{example}

Second, the integral as an \emph{inner product of functions}. In the Fourier and functional-analysis
chapters, the dot product $\vec{u}\cdot\vec{v} = \sum_k u_k v_k$ gets upgraded for functions to
\[
\inner{f}{g} = \int_a^b f(x)\,g(x)\,dx.
\]
The sum over coordinates becomes an integral over the continuum of inputs; everything else --- length
$\norm{f} = \sqrt{\inner{f}{f}}$, orthogonality $\inner{f}{g} = 0$, projection, even Gram--Schmidt
--- transfers verbatim. This is why a chapter on integrals belongs in a linear algebra book at all:
the integral is the continuous analogue of the dot product. Figure~\ref{ch39:fig-innerproduct} makes this bridge
visible: discrete Riemann rectangles melt smoothly into continuous area.

\begin{figure}[H]
\centering
\begin{tikzpicture}[scale=0.95,>=Stealth]
  \begin{scope}[shift={(0,0)}]
    \node[anchor=south,font=\small\bfseries,text=spacegray!90!black] at (2.2,2.7) {(a) Discrete dot product};
    \draw[fill=codebg,draw=spacegray!30,rounded corners=4pt] (-0.4,-0.65) rectangle (4.8,2.5);
    \draw[->,spacegray] (-0.2,0) -- (4.4,0) node[right,font=\scriptsize] {$x$};
    \draw[->,spacegray] (0,-0.2) -- (0,2.3) node[above,font=\scriptsize] {$y$};
    \foreach \x/\h in {0.4/0.6, 1.0/1.2, 1.6/1.7, 2.2/1.85, 2.8/1.5, 3.4/0.9} {
      \draw[fill=vecblue!25,draw=vecblue!70,thick] (\x-0.3,0) rectangle (\x+0.3,\h);
    }
    \node[font=\tiny,vecblue] at (1.6,1.0) {$u_k v_k \Delta x$};
    \node[below,font=\tiny,spacegray] at (0.4,0) {$x_1$};
    \node[below,font=\tiny,spacegray] at (3.4,0) {$x_n$};
    \draw[<->,spacegray] (1.3,-0.22) -- (1.9,-0.22);
    \node[below,font=\tiny,spacegray] at (1.6,-0.22) {$\Delta x$};
  \end{scope}

  \begin{scope}[shift={(6.4,0)}]
    \node[anchor=south,font=\small\bfseries,text=spacegray!90!black] at (2.2,2.7) {(b) Continuous inner product};
    \draw[fill=codebg,draw=spacegray!30,rounded corners=4pt] (-0.4,-0.65) rectangle (4.8,2.5);
    \draw[->,spacegray] (-0.2,0) -- (4.4,0) node[right,font=\scriptsize] {$x$};
    \draw[->,spacegray] (0,-0.2) -- (0,2.3) node[above,font=\scriptsize] {$y$};
    \draw[fill=targetgreen!20,draw=none] (0.1,0) -- plot[domain=0.1:3.7,samples=50] (\x, {1.9*sin(\x*45)}) -- (3.7,0) -- cycle;
    \draw[targetgreen,very thick] plot[domain=0.1:3.7,samples=50] (\x, {1.9*sin(\x*45)});
    \node[targetgreen!80!black,font=\small] at (1.9,0.85) {$\inner{f}{g} = \text{Area}$};
    \node[targetgreen,above right,font=\scriptsize] at (2.4,1.8) {$h(x)=f(x)g(x)$};
    \node[below,font=\tiny,spacegray] at (0.1,0) {$a$};
    \node[below,font=\tiny,spacegray] at (3.7,0) {$b$};
  \end{scope}
\end{tikzpicture}
\caption{The bridge from vector dot products to function inner products. In (a), sampling two
functions at $n$ discrete points turns them into vectors $\vec{u}, \vec{v} \in \R^n$; their weighted
dot product $\sum_{k=1}^n f(x_k)g(x_k)\Delta x$ is the total area of approximating rectangles. In
(b), letting the step size $\Delta x \to 0$ melts the rectangles into the smooth definite integral
$\inner{f}{g} = \int_a^b f(x)g(x)\,dx$. The continuous inner product is simply a Euclidean dot
product evaluated across a continuum of components.}
\label{ch39:fig-innerproduct}
\end{figure}

\begin{connectionbox}
The parallel runs deep. A finite vector is a function from $\{1,\dots,n\}$ to $\R$; a function on
$[a,b]$ is a vector with a continuum of components. Under that dictionary, $\sum_k u_k v_k$ and
$\int u v\,dx$ are the same construction, summation just becomes integration. Orthonormal bases
become Fourier series; the spectral theorem for symmetric matrices becomes the spectral theory of
differential operators. Carrying the geometry of $\R^n$ into spaces of functions is the subject
called \emph{functional analysis}, and the single bridge you cross to get there is replacing $\sum$
by $\int$.
\end{connectionbox}

\begin{focusbox}
The calculus this book needs reduces to four reflexes. (1) A derivative is a slope and a best linear
approximation. (2) The gradient collects the partial derivatives into a vector that points uphill and
vanishes at optima. (3) Setting a derivative or gradient to zero is how optimization meets linear
systems. (4) An integral accumulates, undoes a derivative, and serves as the inner product of
functions. Everything else you can look up.
\end{focusbox}

\subsection*{Limitations and Conventions}

Each of these tools demands respect for its boundaries. Complex numbers guarantee that every
$n\times n$ matrix has a full set of $n$ eigenvalues, but they lose size comparisons:
the complex field $\C$ cannot be ordered, so inequalities like $z_1 < z_2$ have no meaning.
In matrix calculus, literature remains split between numerator layout (gradient as column vector)
and denominator layout (gradient as row vector), causing frequent transpose issues. Finally,
when sums or integrals extend to infinity, linear algebra must lean on analysis: infinite series
$\sum_{k=0}^\infty \mat{A}^k$ require eigenvalue bounds ($\abs{\lambda} < 1$) to converge.

\subsection*{Where We Go Next Across the Curriculum}

Linear algebra connects directly to the rest of the mathematical curriculum:
\begin{itemize}
  \item \emph{Calculus From the Ground Up:} Extends the single-variable slopes and gradients
  met here to multivariable Taylor series, Jacobians, and Hessians.
  \item \emph{Probability and Statistics From the Ground Up:} Interprets quadratic forms
  $\vec{x}\T\mat{\Sigma}\inv\vec{x}$ as Gaussian densities, projections as conditional
  expectations, and SVD as principal component analysis.
  \item \emph{Signal Processing From the Ground Up:} Expands Euler's formula $e^{i\theta}$ and roots
  of unity into the continuous Fourier transform and filter design.
  \item \emph{Algorithms and Optimization From the Ground Up:} Develops gradient steps and normal
  equations into numerical methods for large-scale computation.
\end{itemize}

\section{Exercises}
\label{ch39:exercises}

\begin{exercise}[Unrolling sums]
\label{ch39:ex1}
Evaluate each by writing out the terms.
(a) $\displaystyle\sum_{k=1}^{4}(2k-1)$.\quad
(b) $\displaystyle\sum_{k=0}^{4} \tfrac{1}{2^k}$.\quad
(c) $\displaystyle\sum_{i=1}^{3}\sum_{j=1}^{2} (i + j)$.
For (b), compare your answer with the geometric-series formula~\eqref{ch39:geometric} and with the
infinite sum $\sum_{k=0}^\infty 2^{-k}$.
\end{exercise}

\begin{exercise}[Sigma algebra]
\label{ch39:ex2}
Using only the rules~\eqref{ch39:rule-const}--\eqref{ch39:rule-range} and the closed
forms~\eqref{ch39:sum-ones}--\eqref{ch39:sum-ksq}, show that
\[
\sum_{k=1}^{n}(3k^2 - 2k + 5) = \frac{n(n+1)(2n+1)}{2} - n(n+1) + 5n,
\]
and simplify. Then check your formula at $n=2$ by adding the two terms directly.
\end{exercise}

\begin{exercise}[Order swapping is the trace identity]
\label{ch39:ex3}
For $\mat{A} = \begin{bmatrix}1 & 2\\ 0 & 3\end{bmatrix}$ and $\mat{B} = \begin{bmatrix}4 & 1\\ 1 &
0\end{bmatrix}$, compute $\tr(\mat{A}\mat{B})$ and $\tr(\mat{B}\mat{A})$ from the definition
$\tr(\mat{M}) = \sum_i m_{ii}$, and confirm they are equal. Then explain in one sentence which
summation-order swap makes this true in general.
\end{exercise}

\begin{exercise}[Complex arithmetic]
\label{ch39:ex4}
Let $z = 3 - i$ and $w = 1 + 2i$. Compute $z + w$, $zw$, $\bar z$, $\abs{z}$, and $z/w$ (realize the
denominator). Plot $z$, $w$, and $zw$ in the complex plane and verify by eye that the angle of $zw$
is the sum of the angles of $z$ and $w$.
\end{exercise}

\begin{exercise}[Multiplying is rotating]
\label{ch39:ex5}
Write $z = 2e^{i\pi/3}$ in the form $a + bi$. Then compute $z^3$ two ways: once by cubing $a+bi$
directly, and once with polar form $(re^{i\theta})^3 = r^3 e^{3i\theta}$. Confirm both give the same
number, and describe in words the rotation and scaling that cubing performed.
\end{exercise}

\begin{exercise}[Euler from the series]
\label{ch39:ex6}
Using the power-series argument of Section~\ref{ch39:euler}, evaluate $e^{i\pi/2}$ by separating real
and imaginary parts, and confirm it equals $i$. Then use $e^{i\theta}$ to give a one-line proof of
the double-angle identities $\cos 2\theta = \cos^2\theta - \sin^2\theta$ and $\sin 2\theta =
2\sin\theta\cos\theta$ (hint: expand $(e^{i\theta})^2 = e^{2i\theta}$ two ways).
\end{exercise}

\begin{exercise}[Derivatives from scratch and from rules]
\label{ch39:ex7}
(a) Use the limit definition to show that the derivative of $f(x) = 1/x$ is $-1/x^2$.
(b) Use the product and chain rules to differentiate $g(x) = x^2 e^{-x}$ and find where $g'(x)=0$.
Interpret those points on the graph of $g$.
\end{exercise}

\begin{exercise}[Least squares is calculus]
\label{ch39:ex8}
You have points $(1,1), (2,2), (3,2)$ and want the best-fitting line through the origin, $y = cx$, in
the least-squares sense: minimize $E(c) = \sum_{k}(y_k - c x_k)^2$. Differentiate $E$ with respect to
$c$, set $E'(c) = 0$, and solve for $c$. Then identify which dot products your answer is built from,
and connect it to the formula $c = (\vec{x}\cdot\vec{y})/(\vec{x}\cdot\vec{x})$ for projecting
$\vec{y}$ onto $\vec{x}$.
\end{exercise}

\begin{exercise}[A gradient and its zero]
\label{ch39:ex9}
Let $f(x_1,x_2) = x_1^2 + x_1 x_2 + x_2^2 - 3x_1$. Compute $\nabla f$, set it to the zero vector, and
solve the resulting $2\times 2$ linear system for the critical point. Write that linear system in the
form $\mat{A}\vec{x} = \vec{b}$ and identify the symmetric matrix $\mat{A}$; relate it to the
formula $\nabla f = \mat{A}\vec{x} - \vec{b}$ from Example~\ref{ch39:ex-quadform}.
\end{exercise}

\begin{exercise}[Integration as a dot product]
\label{ch39:ex10}
On the interval $[-\pi, \pi]$ with inner product $\inner{f}{g} = \int_{-\pi}^{\pi} f(x)g(x)\,dx$, show that
$\sin x$ and $\cos x$ are orthogonal, that is $\inner{\sin}{\cos} = 0$. (Hint: $\sin x\cos x =
\tfrac12\sin 2x$.) Then compute $\norm{\sin}^2 = \inner{\sin}{\sin}$. What you have just verified is
the first step in building the Fourier orthonormal basis.
\end{exercise}

\begin{exercise}[Geometric series, scalar to matrix]
\label{ch39:ex11}
(a) Derive $\sum_{k=0}^{n-1} r^k = (1-r^n)/(1-r)$ by the telescoping trick (multiply by $r$ and
subtract). (b) For the matrix $\mat{A} = \begin{bmatrix} 0.5 & 0 \\ 0 & 0.5\end{bmatrix}$, compute
$\mat{I} + \mat{A} + \mat{A}^2 + \cdots$ entry by entry using the scalar formula on the diagonal, and
check that your answer equals $(\mat{I} - \mat{A})\inv$ by inverting $\mat{I}-\mat{A}$ directly.
\end{exercise}

\begin{exercise}[Spiral or decay, from one complex number]
\label{ch39:ex12}
The matrix $\mat{A} = \begin{bmatrix} 0.8 & -0.6\\ 0.6 & 0.8\end{bmatrix}$ has eigenvalues $\lambda =
0.8 \pm 0.6 i$. Compute $\abs{\lambda}$ and $\arg\lambda$ (in degrees). Based on the
Section~\ref{ch39:complex} connection box, predict whether the trajectory $\vec{x}_{k+1} =
\mat{A}\vec{x}_k$ spirals outward, spirals inward, or circles forever, and by what angle per step.
Then scale $\mat{A}$ by $1.25$ and redo the prediction.
\end{exercise}
